\chapter{Control Flow, Functions, Regular Expressions, and Algorithmic Thinking}
\label{chap:control-flow}

\section{Structured Branching Logic}
Execution branching in Python is driven by indentation and short-circuit boolean evaluation.
\begin{lstlisting}
# Short-circuit logic protects against NullPointerException / IndexError
if data is not None and len(data) > 0:
    process_element(data[0])
\end{lstlisting}

\section{Iteration Semantics: Loops and the Exhaustion Invariant}
\begin{definition}[Iteration Protocol]
An object is \textbf{iterable} if it defines \texttt{\_\_iter\_\_()}, which returns an \textbf{iterator}. An iterator defines \texttt{\_\_next\_\_()}, which returns consecutive elements until raising \texttt{StopIteration}.
\end{definition}

\begin{methodbox}[Prime Factorization with Loop Else]
\begin{lstlisting}
def is_prime(n: int) -> bool:
    if n < 2: return False
    for d in range(2, int(n**0.5) + 1):
        if n % d == 0:
            break  # Exits loop early; skips else block
    else:
        return True  # Executed ONLY when range was fully exhausted
    return False
\end{lstlisting}
\end{methodbox}

\section{Function Architecture and the LEGB Scope Rule}
Variable resolution follows the strict LEGB (Local, Enclosing, Global, Built-in) search order.
\begin{figure}[htbp]
\centering
\begin{tcolorbox}[colback=boxnavybg,colframe=brandnavy,title={Scope Resolution Pipeline (LEGB)}]
\small
\textbf{Current Frame (Local)} $\longrightarrow$ \textbf{Outer Functions (Enclosing)} $\longrightarrow$ \textbf{Module Scope (Global)} $\longrightarrow$ \textbf{Python Built-ins}
\end{tcolorbox}
\caption{LEGB Search Resolution Hierarchy}
\end{figure}

\section{Advanced Function Signatures}
\begin{lstlisting}
def configure_pipeline(model_name: str, /, batch_size: int = 32, *args, learning_rate: float = 1e-3, **kwargs):
    '''
    model_name: Positional-only parameter (preceding '/')
    batch_size: Positional or keyword parameter
    *args: Arbitrary positional overflow packed into a tuple
    learning_rate: Keyword-only parameter (succeeding '*')
    **kwargs: Arbitrary keyword overflow packed into a dictionary
    '''
    print(f"Model: {model_name}, Batch: {batch_size}, LR: {learning_rate}")
\end{lstlisting}

\section{Recursion and Memory Stack Frames}
\begin{example}[Recursive Greatest Common Divisor (Euclidean Algorithm)]
\begin{lstlisting}
def gcd(a: int, b: int) -> int:
    if b == 0:
        return a
    return gcd(b, a % b)
\end{lstlisting}
\end{example}
\begin{solutionbox}
\textbf{Trace for \texttt{gcd(48, 18)}:}\\
1. $\operatorname{gcd}(48, 18) \implies b \ne 0 \implies \operatorname{gcd}(18, 48 \pmod{18}) = \operatorname{gcd}(18, 12)$\\
2. $\operatorname{gcd}(18, 12) \implies b \ne 0 \implies \operatorname{gcd}(12, 18 \pmod{12}) = \operatorname{gcd}(12, 6)$\\
3. $\operatorname{gcd}(12, 6) \implies b \ne 0 \implies \operatorname{gcd}(6, 12 \pmod 6) = \operatorname{gcd}(6, 0)$\\
4. $\operatorname{gcd}(6, 0) \implies b = 0 \implies \text{returns } 6$.\\
\textbf{Complexity:} $\mathcal{O}(\log(\min(a, b)))$ steps by Lam\'e's Theorem.
\end{solutionbox}

\section{Pattern Matching with Regular Expressions}
The \texttt{re} engine uses a modified non-deterministic finite automaton (NFA):
\begin{lstlisting}
import re

email_pattern = re.compile(r'^[a-zA-Z0-9_.+-]+@[a-zA-Z0-9-]+\.[a-zA-Z0-9-.]+$')
match = email_pattern.match("student@lpu.co.in")
if match:
    print("Valid academic email address confirmed.")
\end{lstlisting}

\section{Chapter Summary}
\begin{summarybox}
$\bullet$ Functions are first-class citizens: they can be passed as arguments, assigned to variables, and returned from functions.\\
$\bullet$ Default arguments are created once when the function is defined; always use \texttt{None} sentinels for mutable collections.\\
$\bullet$ Pre-compiling regular expressions via \texttt{re.compile()} avoids repetitive string parsing overhead in high-throughput pipelines.
\end{summarybox}
